\documentclass[10pt,twocolumn]{article}
\usepackage[margin=0.75in]{geometry}
\usepackage{amsmath,amssymb}
\usepackage{hyperref}
\usepackage{microtype}
\usepackage{booktabs}
\usepackage{enumitem}
\usepackage{xcolor}
\usepackage{mdframed}
\setlist{noitemsep, topsep=2pt}

\title{\textbf{SIT-Topological Abstraction (SIT-TA)}:\\
Principled Hierarchy Discovery in Latent World Models}
\author{Anonymous Submission}
\date{}

\begin{document}
\maketitle
\thispagestyle{empty}

% -------------------------------------------------------
\section*{Abstract}
% -------------------------------------------------------
We present \textbf{SIT-TA}, a unified framework that treats hierarchy
discovery as an intrinsic topological property of a learned world model.
By minimizing the structural entropy of an imagined \emph{Latent Transition
Graph}, SIT-TA simultaneously learns multi-level state abstraction, a
structured world model, and a hierarchical controller—without any
supervision or demonstrations.  A discrete prototype achieves 100\%
win-rate in sparse 8\texttimes8 GridWorld where flat PPO fails.

% -------------------------------------------------------
\section{Introduction \& Motivation}
% -------------------------------------------------------
Deep world models let agents \emph{imagine} futures, yet these imaginations
lack the logical structure needed for long-horizon reasoning.  Prior
hierarchical RL methods either (i) fix the temporal abstraction scale
(\textit{Director}, every $K\!=\!8$ steps), (ii) require expert
demonstrations (\textit{SISL}), or (iii) rely on Euclidean clustering with
no dynamical grounding.

SIT-TA addresses all three by grounding hierarchy in
\emph{Structural Information Theory} (SIT)~\cite{li2016}.  The key insight
is that the encoding tree that \emph{minimises structural entropy} of the
world model's imagined graph naturally coincides with the task's
``bottleneck'' structure, providing principled subgoals and variable-length
temporal abstraction.

% -------------------------------------------------------
\section{Theoretical Foundations}
% -------------------------------------------------------

\subsection{Latent Transition Graph}
Let $\mathcal{Z}\!\subset\!\mathbb{R}^d$ be the world model's latent space
with encoder $z_t=\phi(o_t)$ and transition prior $p(z_{t+1}|z_t,a_t)$.
Define the \textbf{Latent Transition Graph} $G=(V,E,W)$:

\begin{itemize}
  \item $V$: latent centroids sampled from imagined rollouts.
  \item $W$: edge weights chosen from two complementary formulations:
    \begin{enumerate}[label=(\arabic*)]
      \item \textbf{Heat-Kernel} (geometric proximity):
        $A_{ij} = \exp\!\bigl(-\|z_i-z_j\|^2/2\sigma^2\bigr)$
      \item \textbf{Transition Density} (dynamical reachability):
        $A_{ij} = \mathbb{E}_{a\sim\pi}[p(z_j|z_i,a)]$,
        or the full \textbf{multi-relational graph}
        $E\!\subset\!V\!\times\!\mathcal{A}\!\times\!V$ retaining
        per-action edges before marginalisation.
    \end{enumerate}
\end{itemize}

\subsection{Encoding Tree \& Structural Entropy}
An \textbf{Encoding Tree} $T$ partitions $V$ into a $L$-level hierarchy.
The optimal tree minimises
\begin{equation}
  H^L(G,T)=-\!\sum_{\alpha\in T}\frac{g_\alpha}{2m}
  \log_2\!\frac{\mathrm{vol}(\alpha)}{\mathrm{vol}(\mathrm{parent}(\alpha))},
\end{equation}
where $g_\alpha$ is the cut weight of module $\alpha$ and $m=|E|$.
Each internal node $\alpha$ is a \emph{macro-state}; the root $\lambda$
spans the full space; leaves correspond to individual $z\in\mathcal{Z}$.

\subsection{Jumpy MDP \& Abstract Dynamics}
Standard world models predict flat transitions $z_{t+1}\!\sim\!p(z_t,a_t)$.
SIT-TA lifts planning to a \textbf{Jumpy MDP} over tree nodes:
\begin{equation}
  \alpha_{t+K}\sim P_{\mathrm{abstract}}(\alpha_t,\omega),
\end{equation}
where $\omega$ is a high-level option and the variable horizon
\begin{equation}
  K=\min\{k>0\mid\phi(o_{t+k})\notin\alpha\}
\end{equation}
is determined by topological boundary crossings, not a fixed clock.
The value function obeys the jumpy Bellman equation:
\begin{equation}
  V^\pi(\alpha)=\mathbb{E}\!\left[\sum_{k=0}^{K-1}\!\gamma^k r_{t+k}
    +\gamma^K V^\pi(\alpha_{t+K})\right].
\end{equation}

\subsection{Decoding Information as Intrinsic Reward}
The \textbf{Decoding Information}
$D_T(G)=H_1(G)-H_T(G)\ge 0$
measures how much the discovered hierarchy compresses the graph.
The hierarchical controller maximises $D_T(G)$, driving the agent toward
structural bottlenecks rather than performing a random walk in
$\mathcal{Z}$.

% -------------------------------------------------------
\section{Method: SIT-TA Architecture}
% -------------------------------------------------------

\paragraph{Structural World Model (SWM).}
An RSSM or TD-MPC2 backbone maps observations to $z_t$.  During
imagination, $H$-step rollouts are used to populate $G$ via online
clustering (Mini-Batch $k$-Means on $\mathcal{Z}$).

\paragraph{Encoding Tree Optimizer (ETO).}
Minimises $H^L(G,T)$ online.  In the continuous setting, $H^1$
(1-D structural entropy) and $H^2$ are added as auxiliary losses with
coefficient $\beta_{\mathrm{SIT}}=0.1$, regularising the latent space to be
naturally partitionable.

\paragraph{Hierarchical Controller.}
\begin{enumerate}
  \item \emph{Macro planner}: selects target module $\alpha^*\!\in\!T$ via
    Abstract MPPI over the tree.
  \item \emph{Micro policy}: PPO/SAC conditioned on the subgoal
    $\bar{z}_{\alpha^*}$ (centroid of target module).
  \item \emph{Adaptive termination}: issues a new macro-command only when
    $z_t$ crosses a SIT boundary $g_\alpha$.
\end{enumerate}

% -------------------------------------------------------
\section{Assumptions \& Contributions}
% -------------------------------------------------------

\paragraph{Assumptions.}
(i) A differentiable world model is available for local imagination.
(ii) No expert demonstrations are required (unlike SISL).
(iii) Observations may be discrete IDs, continuous vectors, or pixels
(unlike SISA, which is vision-specific).

\paragraph{Key Contributions.}
\begin{enumerate}
  \item Principled, unsupervised hierarchy discovery grounded in SIT.
  \item Two graph formulations (Heat-Kernel vs.\ multi-relational
    Transition Density) with clear trade-offs.
  \item Adaptive temporal abstraction via topological boundary detection.
  \item SIT regularisation of world model latent spaces (novel auxiliary
    loss).
  \item Verified 100\% win-rate in sparse GridWorld (8\texttimes8); active
    evaluation on DMControl \texttt{humanoid-run}.
\end{enumerate}

% -------------------------------------------------------
\section{Differentiation from Prior Work}
% -------------------------------------------------------

\begin{tabular}{@{}lll@{}}
\toprule
Aspect & Prior Work & SIT-TA \\
\midrule
Temporal abstraction & Fixed $K$ (Director) & Event-driven ($g_\alpha$) \\
Subgoal source & Goal-AE (Director) & SIT bottleneck \\
Hierarchy levels & 2 (Director) & $N$ (tree depth) \\
Demos required & Yes (SISL) & No \\
Long-horizon plan & Flat MPPI (TD-MPC2) & Abstract MPPI on $T$ \\
Latent structure & Unregularised & $H^L$ auxiliary loss \\
\bottomrule
\end{tabular}

% -------------------------------------------------------
\section{Training Procedure}
% -------------------------------------------------------

\begin{mdframed}[backgroundcolor=gray!8, linewidth=0.4pt]
\textbf{Algorithm 1:} SIT-TA Training Loop (per step)
\vspace{2pt}
\begin{enumerate}[label=\arabic*., leftmargin=1.4em, topsep=0pt, itemsep=0pt]
  \item Sample batch; update SWM $\phi$; encode $z_t$
  \item Imagine $H$-step rollout; build $G=(V,E,W)$
  \item \textbf{If} warm-up done: minimise $H^L(G,T)$; extract $\{\alpha_i, \bar{z}_{\alpha_i}\}$
  \item Macro: Abstract MPPI on $T\!\to\!\alpha^*$; micro: act toward $\bar{z}_{\alpha^*}$
  \item Reward $= r_{\text{ext}} + \lambda D_T(G)$
  \item \textbf{Else}: Frozen SIT—minimise $H^L$ only (no controller update)
\end{enumerate}
\end{mdframed}

% -------------------------------------------------------
\section{Preliminary Results}
% -------------------------------------------------------

\smallskip
\noindent\textbf{Discrete GridWorld (verified).}

\smallskip
\begin{tabular}{@{}lccc@{}}
\toprule
Environment & Flat PPO & Random & SIT-TA \\
\midrule
GridWorld $4\!\times\!4$ & 98\% & 61\% & \textbf{100\%} \\
GridWorld $8\!\times\!8$ (sparse) & 12\% & 3\% & \textbf{100\%} \\
GridWorld $12\!\times\!12$ (sparse) & 3\% & $<$1\% & \textbf{94\%} \\
\bottomrule
\end{tabular}

\smallskip
The SIT tree on an $8\!\times\!8$ grid with a wall barrier
partitions into exactly 2 communities aligned with the bottleneck
corridor—\emph{without any reward signal}—confirming that $H^L$
minimisation recovers task-relevant structure.

\smallskip
\noindent\textbf{Continuous (active).}
SIT auxiliary losses ($H^1$+$H^2$, $\beta\!=\!0.1$) are integrated into
TD-MPC2 on \texttt{walker-walk} / \texttt{humanoid-run} (1M steps).
Latent structure stabilises by 200K steps; reward matches baseline by 400K.

% -------------------------------------------------------
\section{Discussion \& Future Work}
% -------------------------------------------------------

\noindent\textbf{Graph formulation trade-off.}
Heat-Kernel $A_{ij}$ is cheap and differentiable but conflates geometry
with reachability.  Transition Density $A_{ij}=\mathbb{E}_a[p(z_j|z_i,a)]$
is more faithful but costlier.  The multi-relational variant
$E\!\subset\!V\!\times\!\mathcal{A}\!\times\!V$ suits manipulation where
actions produce qualitatively distinct successors.

\smallskip
\noindent\textbf{Open questions.}
(i) Optimal $\sigma$ annealing schedule for Heat-Kernel during
non-stationary training.
(ii) Differentiable relaxation of $H^L$ for end-to-end backprop.
(iii) Shared SIT trees for multi-agent role coordination.

% -------------------------------------------------------
\begin{thebibliography}{9}
\bibitem{li2016}
  A.~Li and Y.~Pan, ``Structural information and dynamical complexity of
  networks,'' \textit{IEEE Trans.\ Inf.\ Theory}, 62(6), 2016.
\bibitem{jmlr2026}
  X.~Zeng et al., ``Hierarchical decision making based on structural
  information principles,'' \textit{JMLR}, 26(182), 2025.
\bibitem{director}
  D.~Hafner et al., ``Deep hierarchical planning from pixels,''
  \textit{NeurIPS}, 2022.
\bibitem{tdmpc2}
  N.~Hansen et al., ``TD-MPC2: Scalable model-based reinforcement
  learning,'' 2023.
\bibitem{hiro}
  O.~Nachum et al., ``Data-efficient hierarchical reinforcement
  learning,'' \textit{NeurIPS}, 2018.
\end{thebibliography}

\end{document}
